\chapter{Recurrent Urinary Tract Infections}

\section*{Definition}
Recurrent urinary tract infection (UTI) means repeated microbiologically confirmed infection, conventionally at least two episodes in six months or at least three in one year. Recurrent lower UTIs are common after menopause, but persistent urinary symptoms are not always caused by infection and should be assessed carefully.

\section*{When to See a Doctor}
Seek prompt care for recurrent symptoms, visible blood in the urine, pregnancy, male sex, kidney disease, immunosuppression, or symptoms that do not improve. Urgent assessment is needed for fever, rigors, flank pain, vomiting, confusion, fainting, or severe illness because these may indicate pyelonephritis or sepsis.

\section*{How It Develops}
Reduced estrogen after menopause can thin the urogenital tissues, alter the vaginal microbiome, and reduce local defenses against bacterial colonisation. Infection usually begins with ascent of bowel bacteria through the urethra into the bladder. Repeated antibiotics without culture confirmation can promote resistance and expose patients to avoidable adverse effects.

\section*{Causes}
\begin{itemize}
  \item Menopause-related genitourinary changes, sexual activity, and prior UTI history.
  \item Incomplete bladder emptying, urinary obstruction, pelvic organ prolapse, stones, catheter use, or neurogenic bladder.
  \item Diabetes, immune suppression, pregnancy, and structural urinary-tract abnormalities.
  \item Persistent dysuria from genitourinary syndrome of menopause, vaginitis, urethral irritation, or bladder pain syndrome without infection.
\end{itemize}

\section*{Related Symptoms}
\begin{itemize}
  \item Dysuria, urinary frequency, urgency, suprapubic discomfort, cloudy or strong-smelling urine, and hematuria.
  \item Fever, flank pain, nausea, or vomiting suggest upper-tract infection rather than uncomplicated cystitis.
  \item Vaginal dryness, burning, painful sex, and urgency may reflect genitourinary syndrome of menopause instead of infection.
\end{itemize}

\section*{Medical Evaluation}
\begin{itemize}
  \item Urinalysis and urine culture are used to confirm infection and guide antibiotic selection, particularly with recurrence or atypical symptoms.
  \item Review sexual, gynecologic, menopause, medication, stone, catheter, and bladder-emptying history; examine for prolapse or GSM when indicated.
  \item Imaging or urologic referral is considered for stones, obstruction, hematuria, unusual organisms, relapsing infection, or suspected structural disease.
\end{itemize}

\section*{Treatment Options}
Treat confirmed infection with a culture-directed, appropriately narrow antibiotic course. Repeated recurrences may be prevented with vaginal estrogen when appropriate, behavioural measures, patient-initiated therapy, postcoital prophylaxis, or low-dose continuous prophylaxis after clinician review. Address stones, obstruction, prolapse, retention, diabetes, and other contributing conditions.

\section*{Self-Care and Home Management}
Drink normally to avoid dehydration, do not delay urination habitually, and avoid irritant products. Do not use leftover antibiotics or repeatedly treat urine odor alone. Seek testing when symptoms recur, especially after recent antibiotic exposure.

\section*{References}
\begin{thebibliography}{99}
\bibitem{recurrent_uti:ref1} National Institute for Health and Care Excellence. Urinary tract infection (recurrent): antimicrobial prescribing. NICE guideline NG112; updated 2024.
\bibitem{recurrent_uti:ref2} European Association of Urology. Guidelines on urological infections. EAU; 2025.
\bibitem{recurrent_uti:ref3} The North American Menopause Society. The 2022 hormone therapy position statement. \textit{Menopause}. 2022;29:767--794.
\end{thebibliography}
